\documentclass[11pt,a4paper]{article}

% ---------- Packages ----------
\usepackage[margin=1in]{geometry}
\usepackage{amsmath, amssymb, amsthm, mathtools}
\usepackage{microtype}
\usepackage[colorlinks=true,
            linkcolor=blue!50!black,
            urlcolor=blue!50!black,
            citecolor=blue!50!black]{hyperref}

% ---------- Theorem environments ----------
\theoremstyle{plain}
\newtheorem{theorem}{Theorem}[section]
\newtheorem{proposition}[theorem]{Proposition}
\newtheorem{lemma}[theorem]{Lemma}

\theoremstyle{definition}
\newtheorem{definition}[theorem]{Definition}

\theoremstyle{remark}
\newtheorem{remark}[theorem]{Remark}

% ---------- Math shortcuts ----------
\newcommand{\R}{\mathbb{R}}
\newcommand{\E}{\mathbb{E}}
\newcommand{\F}{\mathcal{F}}
\newcommand{\Prob}{\mathbb{P}}
\newcommand{\Q}{\mathbb{Q}}
\newcommand{\dd}{\mathop{}\!\mathrm{d}}
\newcommand{\eqdef}{\overset{\mathrm{def}}{=}}

% ---------- Document metadata ----------
\title{The Black--Scholes Formula via the Martingale Approach}
\author{Quant Finance Portfolio --- Phase 1, Algorithm 1}
\date{\today}

\begin{document}
\maketitle

\begin{abstract}
We derive the Black--Scholes price of a European call option using the
equivalent martingale measure framework. The argument is self-contained
modulo the standard tools of stochastic calculus---It\^{o}'s lemma,
Girsanov's theorem, and the fundamental theorems of asset pricing---which
we state without proof and apply directly. Beyond the formula itself,
we record three conceptual observations whose appreciation goes beyond
the algebraic manipulation: the distinction between physical and risk-neutral
probabilities of exercise, the change-of-num\'{e}raire interpretation
of \(N(d_1)\), and the way the formula encodes its own replicating
portfolio.
\end{abstract}

% ============================================================
\section{Model}
% ============================================================

Fix a finite horizon \(T > 0\) and a filtered probability space
\((\Omega, \F, \{\F_t\}_{t \in [0,T]}, \Prob)\) supporting a standard
Brownian motion \(W\), with \(\{\F_t\}\) the \(\Prob\)-augmentation of
the natural filtration of \(W\). The market consists of two assets:

\begin{itemize}
    \item A \emph{risky} asset (the stock), whose price \(S\) follows a
    geometric Brownian motion under \(\Prob\):
    \begin{equation}
        \dd S_t = \mu S_t \dd t + \sigma S_t \dd W_t,
        \qquad S_0 > 0,
        \label{eq:gbm-P}
    \end{equation}
    with constants \(\mu \in \R\) and \(\sigma > 0\), paying no dividends.
    \item A \emph{riskless} asset (the money-market account) with price
    \(B_t = e^{rt}\), \(r \geq 0\).
\end{itemize}

We seek the time-\(t\) price \(C_t\) of a European call option with strike
\(K > 0\) and maturity \(T\), whose payoff at expiry is \((S_T - K)^+\).

% ============================================================
\section{Change of Measure}
% ============================================================

Define the discounted stock price \(\tilde{S}_t \eqdef e^{-rt} S_t\).
Applying It\^{o}'s lemma to~\eqref{eq:gbm-P},
\begin{equation}
    \dd \tilde{S}_t
    = (\mu - r) \tilde{S}_t \dd t + \sigma \tilde{S}_t \dd W_t.
    \label{eq:discounted-P}
\end{equation}

For the model to admit no arbitrage, the discounted price must be a
martingale under the pricing measure. The drift in~\eqref{eq:discounted-P}
prevents this under \(\Prob\), so we look for an equivalent measure
\(\Q \sim \Prob\) under which the drift vanishes.

Set the \emph{market price of risk}
\[
    \lambda \eqdef \frac{\mu - r}{\sigma},
\]
and define the Radon--Nikodym derivative
\[
    \left.\frac{\dd \Q}{\dd \Prob}\right|_{\F_T}
    = \exp\!\left(-\lambda W_T - \tfrac{1}{2}\lambda^2 T\right).
\]
Since \(\lambda\) is constant, Novikov's condition
\(\E^{\Prob}\!\left[\exp\!\left(\tfrac{1}{2}\int_0^T \lambda^2 \dd s\right)\right] < \infty\)
holds trivially. By Girsanov's theorem,
\[
    \widetilde{W}_t \eqdef W_t + \lambda t
\]
is a standard Brownian motion under \(\Q\). Substituting
in~\eqref{eq:gbm-P} and~\eqref{eq:discounted-P},
\begin{align}
    \dd S_t
    &= r S_t \dd t + \sigma S_t \dd \widetilde{W}_t,
    \label{eq:gbm-Q} \\
    \dd \tilde{S}_t
    &= \sigma \tilde{S}_t \dd \widetilde{W}_t.
\end{align}
Under \(\Q\), the discounted price is a (local, in fact true) martingale,
and the model is therefore arbitrage-free. Moreover, since the market
has a single source of randomness driving a single risky asset, the
second fundamental theorem of asset pricing implies that \(\Q\) is the
\emph{unique} equivalent martingale measure, and the market is complete.

% ============================================================
\section{The Pricing Formula}
% ============================================================

By the first fundamental theorem of asset pricing, the no-arbitrage
price at time \(t\) of any \(\F_T\)-measurable, sufficiently integrable
contingent claim \(X\) is
\[
    \pi_t(X)
    = B_t \, \E^{\Q}\!\left[\frac{X}{B_T} \,\Big|\, \F_t\right].
\]
For the European call,
\begin{equation}
    C_t
    = e^{-r\tau}
      \E^{\Q}\!\left[(S_T - K)^+ \,\big|\, \F_t\right],
    \qquad \tau \eqdef T - t.
    \label{eq:risk-neutral-pricing}
\end{equation}

% ============================================================
\section{The Distribution of \(S_T\) under \(\Q\)}
% ============================================================

To evaluate~\eqref{eq:risk-neutral-pricing} we need the law of \(S_T\)
under \(\Q\), conditional on \(\F_t\). Applying It\^{o}'s lemma to
\(\log S_t\) using~\eqref{eq:gbm-Q},
\[
    \dd(\log S_t)
    = \left(r - \tfrac{1}{2}\sigma^2\right)\dd t
      + \sigma \dd \widetilde{W}_t.
\]
Integrating from \(t\) to \(T\),
\[
    \log\!\left(\frac{S_T}{S_t}\right)
    = \left(r - \tfrac{1}{2}\sigma^2\right)\tau
      + \sigma\bigl(\widetilde{W}_T - \widetilde{W}_t\bigr),
\]
and since \(\widetilde{W}_T - \widetilde{W}_t \sim \mathcal{N}(0, \tau)\)
under \(\Q\) and is independent of \(\F_t\),
\begin{equation}
    S_T \,\big|\, \F_t
    \;\overset{\Q}{=}\;
    S_t \exp\!\left(
        \left(r - \tfrac{1}{2}\sigma^2\right)\tau
        + \sigma\sqrt{\tau}\, Z
    \right),
    \qquad Z \sim \mathcal{N}(0,1).
    \label{eq:S_T-law}
\end{equation}

% ============================================================
\section{Evaluating the Expectation}
% ============================================================

Let \(\varphi\) and \(N\) denote the standard normal density and CDF
respectively. Substituting~\eqref{eq:S_T-law}
into~\eqref{eq:risk-neutral-pricing},
\[
    C_t
    = e^{-r\tau}
      \int_{\R}\!\left(
          S_t \exp\!\left(
              (r - \tfrac{1}{2}\sigma^2)\tau + \sigma\sqrt{\tau}\, z
          \right) - K
      \right)^+ \!\!\varphi(z) \dd z.
\]
The integrand is positive iff
\[
    S_t \exp\!\left(
        (r - \tfrac{1}{2}\sigma^2)\tau + \sigma\sqrt{\tau}\, z
    \right) > K
    \;\;\Longleftrightarrow\;\;
    z > -d_2,
\]
where
\begin{equation}
    d_2 \eqdef
    \frac{\log(S_t / K) + (r - \tfrac{1}{2}\sigma^2)\tau}
         {\sigma \sqrt{\tau}}.
    \label{eq:d2}
\end{equation}

Splitting the integral over the region where the payoff is positive,
\begin{equation}
    C_t
    = \underbrace{
        e^{-r\tau} S_t
        \int_{-d_2}^{\infty}
            e^{(r - \sigma^2/2)\tau + \sigma\sqrt{\tau}\, z}
            \varphi(z) \dd z
      }_{I_1}
      \;-\;
      \underbrace{
        e^{-r\tau} K
        \int_{-d_2}^{\infty}\!\varphi(z) \dd z
      }_{I_2}.
    \label{eq:split}
\end{equation}

\paragraph{Second integral.} By symmetry of the standard normal density,
\[
    I_2 = e^{-r\tau} K\, \Prob(Z > -d_2)
        = e^{-r\tau} K\, N(d_2).
\]

\paragraph{First integral.} Combine the exponentials in the integrand and
complete the square:
\[
    -\frac{z^2}{2} + \sigma\sqrt{\tau}\, z
    = -\frac{(z - \sigma\sqrt{\tau})^2}{2} + \frac{\sigma^2 \tau}{2}.
\]
The term \(+\sigma^2 \tau / 2\) cancels the \(-\sigma^2 \tau / 2\)
inside \((r - \sigma^2/2)\tau\), and the prefactor \(e^{r\tau}\) cancels
the discount \(e^{-r\tau}\). With the change of variable
\(u \eqdef z - \sigma\sqrt{\tau}\),
\[
    I_1
    = S_t \int_{-d_2 - \sigma\sqrt{\tau}}^{\infty}\!\varphi(u) \dd u
    = S_t\, N\!\left(d_2 + \sigma\sqrt{\tau}\right)
    \eqdef S_t\, N(d_1),
\]
where
\begin{equation}
    d_1 \eqdef d_2 + \sigma\sqrt{\tau}
        = \frac{\log(S_t/K) + (r + \tfrac{1}{2}\sigma^2)\tau}
               {\sigma \sqrt{\tau}}.
    \label{eq:d1}
\end{equation}

% ============================================================
\section{Result}
% ============================================================

Combining the two integrals:

\begin{theorem}[Black--Scholes formula]
The no-arbitrage price at time \(t\) of a European call option with
strike \(K\) and maturity \(T\) on a non-dividend-paying stock
following~\eqref{eq:gbm-P} is
\begin{equation}
    \boxed{\;
        C_t
        = S_t\, N(d_1) \;-\; K e^{-r\tau} N(d_2)
    \;},
    \label{eq:bs}
\end{equation}
where \(\tau = T - t\) and \(d_1, d_2\) are given
by~\eqref{eq:d1}--\eqref{eq:d2}.
\end{theorem}

The price of the corresponding European put follows from put--call
parity, an identity independent of any model assumption beyond
no-arbitrage, which we record for completeness:
\[
    C_t - P_t = S_t - K e^{-r\tau}.
\]
Combined with~\eqref{eq:bs} and the symmetry \(N(x) + N(-x) = 1\), this
gives
\begin{equation}
    P_t = K e^{-r\tau} N(-d_2) - S_t N(-d_1).
    \label{eq:bs-put}
\end{equation}

% ============================================================
\section{Conceptual Remarks}
% ============================================================

\begin{remark}[\(N(d_2)\) is not the physical probability of exercise]
By construction,
\(N(d_2) = \Q(S_T > K \mid \F_t)\), the \emph{risk-neutral} probability
that the option finishes in the money. Under the physical measure, the
same probability would replace \(r\) by \(\mu\) in the numerator of
\(d_2\); the gap between the two is, exactly, the market price of risk
\(\lambda\). Conflating these two probabilities is one of the most
common conceptual errors in introductions to derivatives pricing.
\end{remark}

\begin{remark}[\(N(d_1)\) and the change of num\'{e}raire]
A second probabilistic interpretation, less commonly emphasised, is that
\[
    N(d_1) = \Prob^{S}(S_T > K \mid \F_t),
\]
where \(\Prob^{S}\) is the equivalent measure obtained by taking the
stock \(S\) (rather than the bond \(B\)) as num\'{e}raire. Under
\(\Prob^{S}\), the process \(B_t / S_t\) is a martingale, and the same
change-of-measure machinery used to construct \(\Q\) applies. The
simultaneous appearance of \(N(d_1)\) and \(N(d_2)\) in the formula is
therefore not coincidental: it is the signature of two distinct
num\'{e}raires being used implicitly. This perspective generalises far
beyond Black--Scholes and is the natural setting for pricing exchange
options, LIBOR-based interest-rate products, and most exotic derivatives.
\end{remark}

\begin{remark}[The formula encodes its own hedge]
The decomposition \(C_t = S_t\, N(d_1) - K e^{-r\tau} N(d_2)\) is,
literally, the time-\(t\) value of the replicating portfolio: long
\(N(d_1)\) shares, financed by borrowing \(K e^{-r\tau} N(d_2)\) at the
riskless rate. When in the next algorithm we differentiate~\eqref{eq:bs}
with respect to \(S\) and obtain \(\partial C / \partial S = N(d_1)\)
(the \emph{delta}), with all the \(\varphi\)-terms cancelling cleanly,
this will be a manifestation of the same internal consistency: the
formula is structured around its own hedge.
\end{remark}

\end{document}
